\documentclass[conference]{IEEEtran}
\IEEEoverridecommandlockouts
\usepackage{graphicx}
\usepackage{amsmath}
\usepackage{newtxtext,newtxmath}
\usepackage{url}
\usepackage[hidelinks]{hyperref}
\usepackage{booktabs}
\usepackage{xcolor}
\emergencystretch=3em
\setcounter{dbltopnumber}{2}
\renewcommand{\dbltopfraction}{0.95}
\renewcommand{\dblfloatpagefraction}{0.85}
\renewcommand{\textfraction}{0.05}
\newcommand{\CADDAM}{CADD\allowbreak{}+\allowbreak{}AlphaMissense}
\newcommand{\GCNCADDAM}{GCN\allowbreak{}+\allowbreak{}CADD\allowbreak{}+\allowbreak{}AlphaMissense}

\title{Auditing PPI Graph Augmentation for Missense Variant Pathogenicity Prediction under Topology, Connectivity and Homology Constraints}

\author{Anonymous BIBM submission}

\begin{document}
\maketitle


\begin{abstract}

Protein-protein interaction (PPI) graphs provide a natural source of molecular context for missense variant pathogenicity prediction, but performance gains obtained by propagating protein embeddings over a shared graph can be difficult to interpret. Such gains may reflect protein-topology alignment, train-test connectivity, homology-related split structure, or complementarity to existing variant-effect predictors. We audit these possibilities using a reconstructed ClinVar-STRING dataset of 177,213 filtered, non-conflicting ClinVar missense variants mapped to 14,123 labelled STRING proteins within a 19,622-node STRING protein-association graph used as PPI context. ESM protein language model embeddings are compared with fixed two-hop PPI-GCN propagation, CADD, AlphaMissense and fusion models using validation-selected downstream MLP classifiers. Across three classifier seeds, PPI-GCN increased mean AUC from 0.7835 to 0.7950 and mean AUPRC from 0.6116 to 0.6484. In a retained reference run used for paired graph-ablation analyses, removing direct train-test PPI edges nearly eliminated the AUC gain, protein-clustered bootstrap widened uncertainty across zero, and topology permutation returned performance close to the raw ESM range. Graph fusion also improved \CADDAM{} under the original split, but follow-up controls showed that much of this gain was retained by raw-ESM, edge-ablated and topology-permuted protein representations, indicating limited topology-specific added value in the strong-predictor fusion setting. In the homology-aware stress test, the retained graph-fusion estimate did not exceed the repeated-run architecture-matched \CADDAM{} baseline. Topology-stratified and sample-level transition analyses suggested that graph-associated gains and prediction changes were concentrated in specific network roles and error types. Overall, PPI graph augmentation provided conditional and auditable topological context rather than an unconditional improvement over sequence or variant-effect predictors.

\end{abstract}
\section{Introduction}

Missense variant pathogenicity prediction has been advanced by established methods such as CADD, REVEL, MVP, PrimateAI, EVE and AlphaMissense \cite{Kircher2014CADD,Ioannidis2016REVEL,Qi2021MVP,Sundaram2018PrimateAI,Frazer2021EVE,Cheng2023AlphaMissense}, as well as by protein language models such as ESM \cite{Lin2023ESM}. These approaches capture important sequence, evolutionary and annotation-derived signals. At the same time, disease-associated proteins are embedded in molecular interaction networks, and disease-associated mutations can perturb protein-protein and other macromolecular interactions \cite{Barabasi2011NetworkMedicine,Menche2015DiseaseModules,Sahni2015InteractionPerturbation,Chen2018InteractomePerturbation}. This motivates using PPI graphs as context for missense variant prediction.

Earlier approaches incorporated protein-interaction and other molecular-network features into single-amino-acid variant prioritization, showing that network context can provide information complementary to sequence, structure and annotation features \cite{Yates2014SuSPect,Ozturk2022NetworkFeatures}. More recent graph-based variant predictors include graph attention models for missense effects, heterogeneous gene-variant graph approaches, multimodal variant graph methods and disease-specific graph models \cite{Zhang2022gMVP,Cheng2021VEGN,Yu2025GNNMAP,Ghadie2025DiseaseSpecificGNN}. These studies show that relational structure can support variant interpretation, but they do not resolve a narrower evaluation question: when protein language model embeddings are propagated over a shared STRING PPI graph before variant-level classification, what exactly does the observed gain mean?

This question matters because a shared-graph setting can mix several effects. Train, validation and test variants may have no overlapping protein identifiers, yet their proteins can remain directly connected in the same PPI graph. Related proteins may also cross partitions unless homology is explicitly considered \cite{Littmann2022ProteinMLTips,Teufel2023GraphPart}. Moreover, graph propagation may add little once strong variant-effect predictors such as CADD and AlphaMissense are already present. A rigorous evaluation therefore needs to move beyond a single AUC or AUPRC gain and determine whether the gain is consistent with protein-topology correspondence, sensitive to graph connectivity, robust to protein-level uncertainty, stable under homology-aware evaluation and complementary to established predictors.

We therefore frame the work as an audit rather than a new predictor proposal. Using a reconstructed ClinVar-STRING missense dataset, we evaluate raw ESM representations, PPI-GCN propagation, CADD, AlphaMissense and fusion models. We combine variant-level and protein-clustered bootstrap, train-test edge ablation, partial edge-ablation curves, topology permutation, homology-aware splitting, topology-stratified analyses and corrected-sample profiling. Specifically, we define an audit framework for PPI graph augmentation, test the dependence of raw-ESM-to-PPI-GCN gains on protein-topology correspondence and cross-split connectivity, examine how strong variant-effect predictors and homology-aware splitting bound the apparent gain, and localize graph-associated improvements and errors to specific topology contexts. The central finding is conditional: PPI graph augmentation can improve benchmark scores, but the meaning of the gain depends on the baseline, split design and graph-control setting, and its added value is bounded by homology-aware evaluation and strong external predictors.

\section{Materials and Methods}

\subsection{Dataset and graph construction}

ClinVar variant annotations were obtained from the official NCBI ClinVar FTP tab-delimited variant-summary file \cite{Landrum2020ClinVar}. We used the ClinVar variant-summary release dated 31 March 2024. The exact analysis copy was archived on 3 April 2026, and its checksums are retained in the prepared anonymized source-data archive. We reconstructed a filtered, non-conflicting missense dataset and mapped variants to STRING human protein identifiers. Briefly, we retained GRCh38 germline missense SNVs with criteria-supported review status, combined pathogenic and likely pathogenic records into the positive class, combined benign and likely benign records into the negative class, removed uncertain or conflicting records, and collapsed duplicate VariationIDs. The final ClinVar-STRING dataset contained 177,213 missense variants mapped to 14,123 labelled STRING proteins. The propagation graph contained 19,622 STRING protein nodes and 13,715,404 directed edge-index entries derived from the STRING v12 protein-association network, which was used as a PPI-context graph rather than as a pure physical-interaction map \cite{Szklarczyk2023STRING}.

\subsection{Protein representations and PPI propagation}

Raw protein features were generated using the ESM-2 650M checkpoint, with 1,280-dimensional protein-level vectors obtained by residue averaging \cite{Lin2023ESM}. PPI-GCN embeddings were computed by fixed two-hop GCN-style propagation over the normalized STRING graph \cite{Kipf2017GCN}. We use ``PPI-GCN'' as shorthand for fixed GCN-style propagation of ESM features; no graph-convolution parameters were learned. The implementation symmetrized the retained edge list, added self-loops, coalesced duplicate edges, used symmetric degree normalization and formed the final representation as the average of the raw, one-hop and two-hop features. This propagation step was not trained end-to-end on labels; downstream variant classification used the propagated protein representation corresponding to each variant's mapped STRING protein. Because propagation used the full graph containing train, validation and test nodes and their ESM features, the setting is a transductive shared-graph audit rather than a fully inductive predictor.

\subsection{Splits, classifiers and external predictors}

The original split assigned train, validation and test variants by STRING protein ID, with no shared proteins across partitions. The split contained 141,864 training variants, 17,129 validation variants and 18,220 test variants. The homology-aware split used MMseqs-derived protein clusters so that homologous clusters were assigned wholly to one partition while preserving the same row counts \cite{Steinegger2017MMseqs2}. This split is treated as a stress test rather than a prevalence-matched replacement because the test positive rate changed from 31.95\% in the original split to 10.56\% in the homology-aware split.

All representation and fusion models were evaluated with a downstream MLP classifier selected by validation AUC and optimized with Adam \cite{Kingma2015Adam}. Multi-feature representation/fusion models used a 512/128 hidden-layer MLP with dropout 0.3, whereas low-dimensional external-score baselines used a smaller 32/16 classifier with dropout 0.2 unless explicitly stated otherwise. Training used mini-batches of 1,024 variants, learning rate $10^{-3}$, weight decay $10^{-5}$, a maximum of 50 epochs and early stopping after five consecutive validation checks without improvement in validation AUC. Class imbalance was handled with \texttt{BCEWithLogitsLoss} using training-partition positive-class weight $N_{\mathrm{negative}}/N_{\mathrm{positive}}$. CADD GRCh38 v1.7 scores and AlphaMissense hg38 scores were matched to the reconstructed variants \cite{Kircher2014CADD,Schubach2024CADDv17,Cheng2023AlphaMissense}. CADD scores covered 177,205 variants (99.995\%), and AlphaMissense scores covered 171,753 variants (96.92\%). Missing CADD and AlphaMissense scores were imputed with training-set medians; AlphaMissense-based feature sets also included a binary missingness indicator. In the evaluated MLP configurations, no additional z-score standardization was applied to CADD, AlphaMissense or embedding-derived inputs; the architecture-matched fusion controls therefore share classifier capacity and feature construction, but do not by themselves test alternative feature-scaling choices. Repeated classifier comparisons used seeds 42, 43 and 44 unless otherwise stated.

\subsection{Audit analyses}

Uncertainty was assessed using both variant-level bootstrap and protein-clustered bootstrap \cite{Efron1993Bootstrap}. To test connectivity dependence, we removed direct train-test PPI edges and regenerated PPI-GCN embeddings. Partial edge-ablation curves removed 0\%, 25\%, 50\%, 75\% and 100\% of direct train-test protein pairs using random removal, higher-STRING-score-first removal and lower-STRING-score-first removal. To test whether generic graph smoothing alone could explain the gain, topology permutation preserved the unlabeled graph topology while randomizing the correspondence between protein identity and graph position. Follow-up fusion controls also concatenated \CADDAM{} features with raw ESM, edge-ablated PPI-GCN and topology-permuted PPI-GCN representations to distinguish protein-representation complementarity from graph-topology-specific complementarity.

Topology-context analyses stratified test variants by degree, k-core, betweenness, clustering coefficient, neighbour-risk and community-boundary features. Neighbour-risk and neighbour pathogenic-rate features for validation and test proteins were computed from training-set labels only; validation and test labels were not used to construct these topology-context features. Corrected-sample analyses compared variants whose predictions were corrected or degraded after adding graph context. Multiple topology-feature comparisons used Benjamini-Hochberg false-discovery-rate correction \cite{Benjamini1995FDR}. AUPRC was interpreted within each split because precision-recall metrics are sensitive to class prevalence \cite{Saito2015PrecisionRecall}.

\section{Results}

\subsection{Full-graph PPI-GCN improves point estimates but depends on graph connectivity}

Under the original protein-level split, a retained reference run used for paired graph-ablation and bootstrap analyses gave raw ESM test AUC 0.7790 and AUPRC 0.6074. PPI-GCN improved these values to AUC 0.7997 and AUPRC 0.6496, corresponding to $\Delta$ AUC = +0.0207 and $\Delta$ AUPRC = +0.0421 (Fig. 1B). Across three repeated classifier seeds, the comparison remained directionally positive but with a smaller mean AUC difference: raw ESM reached AUC 0.7835 $\pm$ 0.0051 and AUPRC 0.6116 $\pm$ 0.0054, whereas PPI-GCN reached AUC 0.7950 $\pm$ 0.0013 and AUPRC 0.6484 $\pm$ 0.0096. Variant-level bootstrap of the retained reference run supported a positive difference, with $\Delta$ AUC +0.0208, 95\% CI [0.0174, 0.0240], and $\Delta$ AUPRC +0.0420, 95\% CI [0.0353, 0.0484]. However, protein-clustered bootstrap widened uncertainty across zero: $\Delta$ AUC CI [-0.0080, +0.0527] and $\Delta$ AUPRC CI [-0.0021, +0.0798] (Fig. 1C).

The full graph contained 1,242,422 direct train-test edge-index entries, representing 9.06\% of all edge-index entries and 621,211 unique undirected train-test protein pairs. Removing all direct train-test PPI edges reduced edge-ablated ESM-GCN to AUC 0.7791 $\pm$ 0.0076 and AUPRC 0.6255 $\pm$ 0.0148 (Fig. 1D). Relative to raw ESM, the AUC gain was nearly eliminated ($\Delta$ AUC +0.0001), while a smaller AUPRC gain remained ($\Delta$ AUPRC +0.0181). Partial ablation further showed that AUC gain declined as direct train-test pairs were removed, with higher-STRING-score edge removal generally causing stronger loss at intermediate and high removal fractions (Fig. 1E).

Topology permutation provided a second boundary. When the graph topology was preserved but protein identities were randomly reassigned to graph positions, topology-permuted ESM-GCN achieved AUC 0.7757 $\pm$ 0.0067 and AUPRC 0.6052 $\pm$ 0.0086, close to raw ESM and below real PPI-GCN (Fig. 1F). Thus, the retained full-graph gain was consistent with protein-specific topology alignment rather than generic graph smoothing, although its AUC component was also strongly sensitive to cross-split connectivity.

\begin{figure*}[!t]
\centering
\includegraphics[width=0.98\textwidth]{figure_drafts/Figure1_connectivity_topology_controls_v2_accepted.pdf}
\caption{PPI graph augmentation is sensitive to protein-topology correspondence and train-test connectivity. A, Original protein-level split on a shared PPI graph. B, Retained reference run in which full-graph PPI-GCN improves over raw ESM; repeated classifier seeds show the same direction with a smaller mean AUC difference, as reported in the text and Supplementary Table 2. C, Variant-level bootstrap supports positive gain, whereas protein-clustered bootstrap widens uncertainty across zero. D, Removing direct train-test PPI edges nearly eliminates the reference-run AUC gain. E, Partial train-test edge ablation shows connectivity-sensitive degradation. F, Topology permutation returns performance close to raw ESM.}
\label{fig:connectivity}
\end{figure*}

\subsection{PPI context is complementary under the original split but bounded by strong predictors and homology-aware evaluation}

Under the original split, the external variant-effect predictors provided strong baselines. The single-score CADD and AlphaMissense baselines achieved CADD-only AUC 0.9207 and AUPRC 0.8049, and AlphaMissense-only AUC 0.9470 and AUPRC 0.8899. In the matched repeated-run fusion-control audit used for Fig. 2, an architecture-matched \CADDAM{} baseline trained with the same 512/128 MLP used for fusion controls achieved mean AUC 0.9572 and AUPRC 0.9068. Adding raw ESM representation to \CADDAM{} improved performance to AUC 0.9680 and AUPRC 0.9317, indicating that protein-level sequence representation itself was complementary to the two scalar predictors. In the matched repeated-run graph-fusion comparison, real PPI-GCN fusion achieved AUC 0.9718 and AUPRC 0.9415, but edge-ablated graph fusion remained close at AUC 0.9708 and AUPRC 0.9395, and topology-permuted graph fusion retained AUC 0.9694 and AUPRC 0.9355. Paired bootstrap using ensemble-averaged predictions supported a real-fusion margin over raw ESM fusion (protein-clustered $\Delta$ AUC 95\% CI 0.0005 to 0.0075; $\Delta$ AUPRC CI 0.0032 to 0.0167), but the smaller margins over edge-ablated and topology-permuted graph controls crossed zero under protein-clustered uncertainty. Thus, in the strong-predictor fusion setting, the original-split gain was largely retained by protein-representation controls, while the topology-specific margin of real PPI-GCN over graph controls was small and not consistently supported under clustered uncertainty (Fig. 2B,C).

The homology-aware split changed the interpretation. PPI-GCN remained higher than raw ESM, increasing AUC from 0.7248 to 0.7606 and AUPRC from 0.2418 to 0.2768 (Fig. 2D). However, once CADD and AlphaMissense were included, the retained graph-fusion estimate did not exceed the repeated-run architecture-matched two-predictor baseline. The architecture-matched \CADDAM{} baseline achieved AUC 0.9754 and AUPRC 0.8401; \CADDAM{} plus raw ESM achieved mean AUC 0.9728 and AUPRC 0.8249; and the retained graph-fusion model achieved AUC 0.9741 and AUPRC 0.8333 (Fig. 2E). An earlier retained-reference bootstrap likewise did not support a positive margin, although this was not a fully matched repeated-ensemble comparison. These results show that PPI context can be useful relative to raw ESM, but its marginal value is bounded when strong variant-effect predictors and stricter split structure are considered.

\begin{figure*}[!t]
\centering
\includegraphics[width=0.98\textwidth]{figure_drafts/Figure2_fusion_homology_v1_accepted.pdf}
\caption{PPI graph context is bounded by protein-representation controls and homology-aware evaluation. A, CADD and AlphaMissense matching and fusion-control setup. B, Original-split repeated-run performance of architecture-matched \CADDAM{}, \CADDAM{}+raw ESM, \CADDAM{}+edge-ablated PPI-GCN, \CADDAM{}+topology-permuted PPI-GCN and \CADDAM{}+real PPI-GCN. Bars show mean test performance across three matched classifier seeds except topology-permuted fusion, which uses three topology permutations by three classifier seeds (n=9); error bars indicate one standard deviation. C, Marginal test-set gain of real PPI-GCN fusion over architecture-matched \CADDAM{} and protein-representation controls; error bars show protein-clustered paired-bootstrap 95\% confidence intervals computed from ensemble-averaged test predictions. Because panel C computes metrics from ensemble-averaged probabilities, its point deltas need not equal differences between the across-seed mean metrics shown in panel B. The topology-permuted interval conditions on the sampled topology permutations and does not include permutation-generation uncertainty. D, Homology-aware split comparison of raw ESM and PPI-GCN. E, Homology-aware architecture-matched \CADDAM{}, \CADDAM{}+raw ESM and graph fusion; \CADDAM{} and \CADDAM{}+raw ESM are repeated-run means, whereas graph fusion is the retained homology-aware point estimate. F, Test-set prevalence differs between original and homology-aware splits.}
\label{fig:fusion_homology}
\end{figure*}

\subsection{Graph gains concentrate in specific topology contexts}

ESM representation space showed weak but detectable correspondence with PPI topology. Sampled PPI-neighbour pairs had mean ESM cosine similarity 0.8789 compared with 0.8639 for random protein pairs, and within-community pairs had higher similarity than between-community pairs (0.8743 versus 0.8642; Fig. 3A). This pattern is consistent with modest correspondence between ESM representation space and PPI network structure before propagation.

Graph-derived gains were not monotonic with simple centrality. Mid-k-core variants showed $\Delta$ AUC +0.0775 and $\Delta$ AUPRC +0.1049, whereas high-degree variants showed $\Delta$ AUC -0.0178 (Fig. 3B). Betweenness and clustering also shaped gain profiles: low- and mid-betweenness groups showed positive AUC gains (+0.0614 and +0.0522), while high betweenness showed a negative AUC difference (-0.0145); mid- and high-clustering groups showed positive gains, while low clustering was close to zero (Fig. 3C).

Neighbour-risk stratification and joint topology maps showed that the graph signal was concentrated in local contexts rather than uniformly distributed. Low-, mid- and high-neighbour-risk strata all had positive AUC gains (+0.0641, +0.0265 and +0.0284), but the shape of AUC and AUPRC gains differed across strata (Fig. 3D). Exploratory joint topology maps further identified stronger positive and negative profiles, including mid-k-core plus high-betweenness variants ($\Delta$ AUC +0.1429), mid-k-core plus mid-clustering variants (+0.1408), high-neighbour-risk plus mid-k-core variants (+0.1274), and high-k-core plus low-clustering variants (-0.0532; Fig. 3E,F). These cells are descriptive predictive gain profiles selected within this dataset, not independently validated biological subtypes; their bootstrap intervals quantify within-cell sampling uncertainty rather than selection-adjusted reproducibility.

\begin{figure*}[!t]
\centering
\includegraphics[width=0.98\textwidth]{figure_drafts/Figure3_topology_context_v1_accepted.pdf}
\caption{PPI-GCN gains concentrate in specific network topology contexts. A, ESM representations show weak alignment with PPI topology. B, Degree and k-core gains are non-monotonic. C, Betweenness and clustering strata shape graph-derived gain. D, Neighbour-risk strata show local-context dependence. E, Joint topology cells contain both positive and negative gains. F, Selected exploratory joint-cell profiles illustrate heterogeneous gains. The joint-cell panels are descriptive and were not used for cell-level multiple-testing inference.}
\label{fig:topology_context}
\end{figure*}

\subsection{Graph augmentation reshapes false-negative and false-positive predictions differently}

At a fixed decision-score threshold of 0.5, raw ESM and PPI-GCN were jointly correct for 11,875 test variants and jointly wrong for 3,965 variants. PPI-GCN corrected 1,436 variants and degraded 944 variants relative to raw ESM. Corrected variants included 837 pathogenic false negatives rescued and 599 benign false positives suppressed, whereas degraded variants included 904 benign true negatives converted to false positives and 40 pathogenic true positives converted to false negatives (Fig. 4A).

As an exploratory retained-reference analysis, the transition from the earlier \CADDAM{} baseline to the corresponding retained graph-fusion model had a different structure from the raw-ESM transition. The two retained models were jointly correct for 15,981 variants and jointly wrong for 1,298 variants; graph fusion corrected 666 variants and degraded 275 variants. The corrected set contained more suppressed false positives than rescued false negatives (444 versus 222; Fig. 4B). This transition analysis was not recomputed from the architecture-matched repeated-run ensembles used in Fig. 2.

Topology profiles differed between rescued false negatives and suppressed false positives. In the raw ESM to PPI-GCN comparison, rescued false negatives had higher degree, neighbour-risk, neighbour pathogenic rate, neighbour training-variant count and betweenness, whereas suppressed false positives had higher clustering coefficient and within-community neighbour fraction (Fig. 4C). Similar but not identical profile shifts appeared in the retained \CADDAM{} to graph-fusion comparison (Fig. 4D). Multiple topology features remained significant after FDR correction (Fig. 4E), and difficulty-matched contrasts retained differences across degree, k-core, neighbour-risk, neighbour pathogenic rate, betweenness and boundary role (Fig. 4F). These results associate graph-related prediction changes with network context while remaining descriptive rather than causal.

\begin{figure*}[!t]
\centering
\includegraphics[width=0.98\textwidth]{figure_drafts/Figure4_prediction_transitions_v2_accepted.pdf}
\caption{Graph augmentation changes false-negative and false-positive predictions in different network contexts. A, Prediction transitions from raw ESM to PPI-GCN at a fixed decision-score threshold of 0.5. B, Retained-reference \CADDAM{} to graph-fusion transitions; this exploratory transition analysis was not recomputed from the architecture-matched repeated-run ensembles in Fig. 2. C-D, Row-normalized heatmaps compare topology profiles of suppressed false positives and rescued false negatives. E, Corrected groups differ across multiple topology features after FDR correction. F, Difficulty-matched contrasts retain topology differences.}
\label{fig:prediction_transitions}
\end{figure*}

\section{Discussion}

This audit shows that PPI graph augmentation can improve benchmark performance in a shared-graph evaluation setting, but the meaning of the gain depends on evaluation design and baseline choice. In the original shared-graph split, PPI-GCN improved over raw ESM, and this raw-representation gain was sensitive to protein-topology correspondence and cross-split connectivity. In contrast, once \CADDAM{} was included, follow-up controls showed that adding raw ESM, edge-ablated PPI-GCN or topology-permuted PPI-GCN all produced large fusion gains, with real PPI-GCN only modestly higher. Accordingly, the strong-predictor fusion gain should not be interpreted as purely PPI-topology-specific; it reflects substantial protein-representation complementarity plus a smaller real-topology margin.

The homology-aware analysis adds a second boundary. PPI-GCN remained useful relative to raw ESM, but the retained graph-fusion estimate did not exceed the repeated-run architecture-matched \CADDAM{} baseline once homologous protein clusters were kept within partitions, and available retained-reference bootstrap evidence did not support a positive margin. Because the homology-aware split also changed test-set prevalence, it should be read as a stress test rather than a prevalence-matched replacement. Nevertheless, the result reinforces the need to condition claims about graph complementarity on the baseline predictor, split design and graph-control baseline.

The topology and transition analyses suggest that graph propagation changes predictions in structured ways. Gains concentrated in specific network roles, and rescued false negatives differed from suppressed false positives across multiple topology features. These patterns are useful for understanding model behaviour, but they do not prove a causal biological mechanism. Future work should test whether such profiles replicate in independent cohorts, temporal ClinVar splits, disease-specific networks and alternative PPI sources.

Several limitations and threats to validity remain. ClinVar labels may reflect submission bias and future label updates \cite{Landrum2020ClinVar}. We did not explicitly audit historical overlap between the reconstructed ClinVar benchmark and variants potentially used during development, calibration or prior evaluation of external predictors such as CADD and AlphaMissense; their absolute performance should therefore be interpreted as benchmark performance rather than as a fully independent prospective estimate. STRING combined scores integrate multiple evidence channels and are not equivalent to direct physical interaction strength \cite{Szklarczyk2023STRING}. The original split is a transductive shared full-graph setting rather than a fully disconnected inductive evaluation. Removing direct train-test edges does not remove train-validation, validation-test or higher-order paths, so edge ablation should be interpreted as a targeted connectivity test rather than a complete leakage removal procedure. Validation AUC was used for checkpoint selection, and the retained reference PPI-GCN run showed a larger validation-set AUC gain than test-set gain (validation AUC 0.8237 versus 0.7814 for raw ESM; test AUC 0.7997 versus 0.7790). The main conclusions therefore rely on held-out test metrics, repeated-seed summaries and protein-clustered uncertainty rather than validation gains alone. The homology-aware split reduces homologous-cluster sharing but changes class prevalence, which limits direct AUPRC comparisons across splits. The transition counts in Fig. 4 depend on a fixed decision-score threshold, and class-weighted sigmoid outputs should not be interpreted as calibrated clinical probabilities. Alternative feature-scaling strategies were not evaluated, so the magnitude of representation-complementarity gains may depend partly on the retained input-scaling and optimization setup. Finally, this audit relies on retained processed outputs and a compact source-data archive; large provider-controlled resources such as CADD and AlphaMissense must be obtained from their original sources. Taken together, these limitations support treating PPI graph augmentation as a conditional source of network context rather than as a universally reliable substitute for strong sequence or variant-effect predictors.

\section*{Data and Code Availability}

A compact anonymized source-data archive containing figure source data, summary metrics, bootstrap outputs and checksums has been prepared. The archive and analysis scripts will be released according to venue policy. Large third-party resources, including CADD, AlphaMissense, ClinVar and STRING source files, should be obtained from their original providers.

\section*{Ethical Statement and AI-Tool Disclosure}

This study used publicly available aggregate variant annotations and computationally derived representations. It did not involve direct recruitment of human participants, intervention, or access to identifiable individual-level clinical records.

OpenAI ChatGPT, OpenAI Codex, and Anthropic Claude were used for language editing, literature organization, consistency checking, and assistance with manuscript drafting, figure-caption preparation, and analysis-script drafting and debugging. Study design, model comparisons, evaluation criteria, experiment execution, result interpretation, numerical verification, and citation checking remained under human control. All reported numerical results and citations were verified against retained analysis outputs, source-data files, and original references.

The AI tools did not act as authors, independently determine the scientific conclusions, or generate or fabricate experimental data. Full responsibility for the study and the final manuscript rests with the authors.



\begin{thebibliography}{99}
\bibitem{Barabasi2011NetworkMedicine} A.-L. Barabasi, N. Gulbahce, and J. Loscalzo, ``Network medicine: a network-based approach to human disease,'' \emph{Nature Reviews Genetics}, vol. 12, pp. 56--68, 2011, doi: 10.1038/nrg2918.
\bibitem{Benjamini1995FDR} Y. Benjamini and Y. Hochberg, ``Controlling the false discovery rate: a practical and powerful approach to multiple testing,'' \emph{Journal of the Royal Statistical Society: Series B}, vol. 57, pp. 289--300, 1995, doi: 10.1111/j.2517-6161.1995.tb02031.x.
\bibitem{Chen2018InteractomePerturbation} S. Chen et al., ``An interactome perturbation framework prioritizes damaging missense mutations for developmental disorders,'' \emph{Nature Genetics}, vol. 50, pp. 1032--1040, 2018, doi: 10.1038/s41588-018-0130-z.
\bibitem{Cheng2021VEGN} J. Cheng, C. Lawrence, and M. Niepert, ``VEGN: variant effect prediction with graph neural networks,'' arXiv:2106.13642, 2021, doi: 10.48550/arXiv.2106.13642.
\bibitem{Cheng2023AlphaMissense} J. Cheng et al., ``Accurate proteome-wide missense variant effect prediction with AlphaMissense,'' \emph{Science}, vol. 381, eadg7492, 2023, doi: 10.1126/science.adg7492.
\bibitem{Efron1993Bootstrap} B. Efron and R. J. Tibshirani, \emph{An Introduction to the Bootstrap}. Chapman \& Hall/CRC, 1993.
\bibitem{Frazer2021EVE} J. Frazer et al., ``Disease variant prediction with deep generative models of evolutionary data,'' \emph{Nature}, vol. 599, pp. 91--95, 2021, doi: 10.1038/s41586-021-04043-8.
\bibitem{Ghadie2025DiseaseSpecificGNN} M. Ghadie, S. Sardaar, and Y. Trakadis, ``Disease-specific prediction of missense variant pathogenicity with DNA language models and graph neural networks,'' \emph{Bioengineering}, vol. 12, p. 1098, 2025, doi: 10.3390/bioengineering12101098.
\bibitem{Ioannidis2016REVEL} N. M. Ioannidis et al., ``REVEL: an ensemble method for predicting the pathogenicity of rare missense variants,'' \emph{American Journal of Human Genetics}, vol. 99, pp. 877--885, 2016, doi: 10.1016/j.ajhg.2016.08.016.
\bibitem{Kingma2015Adam} D. P. Kingma and J. Ba, ``Adam: a method for stochastic optimization,'' \emph{International Conference on Learning Representations}, 2015, arXiv:1412.6980.
\bibitem{Kipf2017GCN} T. N. Kipf and M. Welling, ``Semi-supervised classification with graph convolutional networks,'' \emph{International Conference on Learning Representations}, 2017.
\bibitem{Kircher2014CADD} M. Kircher et al., ``A general framework for estimating the relative pathogenicity of human genetic variants,'' \emph{Nature Genetics}, vol. 46, pp. 310--315, 2014, doi: 10.1038/ng.2892.
\bibitem{Landrum2020ClinVar} M. J. Landrum et al., ``ClinVar: improvements to accessing data,'' \emph{Nucleic Acids Research}, vol. 48, pp. D835--D844, 2020, doi: 10.1093/nar/gkz972.
\bibitem{Lin2023ESM} Z. Lin et al., ``Evolutionary-scale prediction of atomic-level protein structure with a language model,'' \emph{Science}, vol. 379, pp. 1123--1130, 2023, doi: 10.1126/science.ade2574.
\bibitem{Littmann2022ProteinMLTips} Q. Hou, K. Waury, D. Gogishvili, and K. A. Feenstra, ``Ten quick tips for sequence-based prediction of protein properties using machine learning,'' \emph{PLOS Computational Biology}, vol. 18, no. 12, Art. no. e1010669, 2022, doi: 10.1371/journal.pcbi.1010669.
\bibitem{Menche2015DiseaseModules} J. Menche et al., ``Disease networks. Uncovering disease-disease relationships through the incomplete interactome,'' \emph{Science}, vol. 347, 1257601, 2015, doi: 10.1126/science.1257601.
\bibitem{Ozturk2022NetworkFeatures} K. Ozturk and H. Carter, ``Predicting functional consequences of mutations using molecular interaction network features,'' \emph{Human Genetics}, vol. 141, pp. 1195--1210, 2022, doi: 10.1007/s00439-021-02329-5.
\bibitem{Qi2021MVP} H. Qi et al., ``MVP predicts the pathogenicity of missense variants by deep learning,'' \emph{Nature Communications}, vol. 12, p. 510, 2021, doi: 10.1038/s41467-020-20847-0.
\bibitem{Sahni2015InteractionPerturbation} N. Sahni et al., ``Widespread macromolecular interaction perturbations in human genetic disorders,'' \emph{Cell}, vol. 161, pp. 647--660, 2015, doi: 10.1016/j.cell.2015.04.013.
\bibitem{Saito2015PrecisionRecall} T. Saito and M. Rehmsmeier, ``The precision-recall plot is more informative than the ROC plot when evaluating binary classifiers on imbalanced datasets,'' \emph{PLOS ONE}, vol. 10, e0118432, 2015, doi: 10.1371/journal.pone.0118432.
\bibitem{Schubach2024CADDv17} M. Schubach et al., ``CADD v1.7: using protein language models, regulatory CNNs and other nucleotide-level scores to improve genome-wide variant predictions,'' \emph{Nucleic Acids Research}, vol. 52, no. D1, pp. D1143--D1154, 2024, doi: 10.1093/nar/gkad989.
\bibitem{Steinegger2017MMseqs2} M. Steinegger and J. Soding, ``MMseqs2 enables sensitive protein sequence searching for the analysis of massive data sets,'' \emph{Nature Biotechnology}, vol. 35, pp. 1026--1028, 2017, doi: 10.1038/nbt.3988.
\bibitem{Sundaram2018PrimateAI} L. Sundaram et al., ``Predicting the clinical impact of human mutation with deep neural networks,'' \emph{Nature Genetics}, vol. 50, pp. 1161--1170, 2018, doi: 10.1038/s41588-018-0167-z.
\bibitem{Szklarczyk2023STRING} D. Szklarczyk et al., ``The STRING database in 2023: protein-protein association networks and functional enrichment analyses for any sequenced genome of interest,'' \emph{Nucleic Acids Research}, vol. 51, pp. D638--D646, 2023, doi: 10.1093/nar/gkac1000.
\bibitem{Teufel2023GraphPart} F. Teufel, M. H. Gislason, J. J. Almagro Armenteros, A. R. Johansen, O. Winther, and H. Nielsen, ``GraphPart: homology partitioning for biological sequence analysis,'' \emph{NAR Genomics and Bioinformatics}, vol. 5, lqad088, 2023, doi: 10.1093/nargab/lqad088.
\bibitem{Yates2014SuSPect} C. M. Yates, I. Filippis, L. A. Kelley, and M. J. E. Sternberg, ``SuSPect: enhanced prediction of single amino acid variant phenotype using network features,'' \emph{Journal of Molecular Biology}, vol. 426, pp. 2692--2701, 2014, doi: 10.1016/j.jmb.2014.04.026.
\bibitem{Yu2025GNNMAP} H. Yu et al., ``A graph neural network approach for accurate prediction of pathogenicity in multi-type variants,'' \emph{Briefings in Bioinformatics}, vol. 26, bbaf151, 2025, doi: 10.1093/bib/bbaf151.
\bibitem{Zhang2022gMVP} H. Zhang, M. S. Xu, X. Fan, W. K. Chung, and Y. Shen, ``Predicting functional effect of missense variants using graph attention neural networks,'' \emph{Nature Machine Intelligence}, vol. 4, no. 11, pp. 1017--1028, 2022, doi: 10.1038/s42256-022-00561-w.
\end{thebibliography}
\end{document}
